% explainer-source-hash: f749393b7f38e0d6d15c3f1e7bc1a051ad8c17439c55d82b1a6fe2bb4aefe790
% explainer-source-commit: a49871e563da2cb89c2448fb6db326242aa5301c
% explainer-generated: 2026-07-16
% Regenerate with /explain-all (see WIRING.md). Do not edit this stamp by hand:
% it is how the toolchain knows whether this document still matches the code.
\documentclass[11pt,a4paper]{article}
\input{preamble}

\title{\vspace{-1.5cm}\textbf{Paws Rescue Donations}\\[0.3cm]
\large A Deep Dive into a Django Donation Tracker}
\author{Explainer for \texttt{paws-rescue-donations}}
\date{}

\begin{document}
\maketitle
\thispagestyle{empty}

\begin{keyidea}{What you are about to read}
This document explains every meaningful part of the \file{paws-rescue-donations}
codebase, assuming you have never written a line of Python or seen a web application
before. Every technical term is defined the first time it appears. The code, not the
README, is the authority throughout: where the two disagree, that disagreement is
reported as a finding rather than smoothed over.
\end{keyidea}

\tableofcontents
\newpage

\section{The project in one page}

\subsection{The problem it solves}

An animal rescue takes in money. Somebody donates 1000 rupees and earmarks it for a
specific cat named Bella. Somebody else donates 500 for a dog named Max. A volunteer
needs to answer, later, questions like ``how much has Bella received this month'' and
``show me every bank transfer above 2000 rupees''.

The usual tool for this is a spreadsheet. A spreadsheet works until several people are
entering records at once, at which point three problems arrive together: there is no
access control (anyone with the link can edit or delete anything), there are no
consistent fields (one person types ``Bank'', another types ``bank transfer''), and
there is no reliable way to ask a question of the data.

This project replaces the spreadsheet with a small database-backed web application
sitting behind a login. It was built for a faculty-supervised Khidmat community-service
project at Habib University.

\begin{jargon}{Web application}
A program that runs on a computer somewhere else (the \emph{server}) and that you use
through a web browser. The browser sends a \emph{request} (``show me the dashboard''),
the server does some work and sends back a \emph{response} (a page of HTML that the
browser draws on screen).
\end{jargon}

\begin{jargon}{Database}
A program dedicated to storing structured records and answering questions about them
quickly and safely, even when many users read and write at the same time. This project
uses SQLite for local development and PostgreSQL in production. Both are
\emph{relational} databases: they store data in tables of rows and columns, like a
spreadsheet, but with enforced rules about what may go in each column.
\end{jargon}

\subsection{The shape of the whole thing}

The application has exactly one kind of user: a staff member who logs in. There is no
public-facing side at all, no donor-visible page, no sign-up form. Everything is behind
the login.

\begin{keyidea}{The entire feature set}
\begin{itemize}
  \item Log in with an email address and password.
  \item Create, read, update and delete donation records.
  \item Filter and sort those records seven different ways.
  \item See running totals and a ``top 5 pets'' panel on a dashboard.
  \item Export the current filtered view as a CSV file.
  \item Add and remove other staff logins from the web interface.
\end{itemize}
\end{keyidea}

\subsection{Scale and honesty about it}

This is a small project, and this document treats it as one. The entire application is
roughly 900 lines of Python plus about 600 lines of HTML templates. There is no
background job queue, no cache, no message broker, no microservice, and no JavaScript
written by hand. That is a reasonable set of choices for the problem, and the document
does not inflate them into an architecture they are not.

\section{Django, briefly, for someone who has never met it}

Every file in this project is arranged the way it is because Django expects it that way.
Understanding the framework's vocabulary makes the rest of the document read easily.

\begin{jargon}{Framework}
A pre-built skeleton of a program. Rather than writing the plumbing that listens for
browser requests, matches URLs, talks to the database and renders HTML, you write only
the parts specific to your problem and drop them into slots the framework provides.
Django is a web framework for Python.
\end{jargon}

\begin{jargon}{MVT (Model, View, Template)}
Django's way of splitting a web application into three jobs.
\textbf{Model}: a Python class describing one kind of record, which Django maps onto a
database table. \textbf{View}: a Python function that receives a request, decides what
to do, and returns a response. \textbf{Template}: an HTML file with placeholders that
get filled in with data. Django calls this MVT; other frameworks call the same idea MVC.
\end{jargon}

\begin{jargon}{ORM (Object-Relational Mapper)}
The layer that lets you write \code{Donation.objects.filter(pet\_name='Max')} in Python
instead of writing \code{SELECT * FROM donations WHERE pet\_name = 'Max'} in SQL. It
translates Python method calls into database queries and turns the returned rows back
into Python objects.
\end{jargon}

\begin{jargon}{QuerySet}
The object the ORM hands back when you ask for records. Its defining property is
\emph{laziness}: it does not touch the database when you build it, only when you
actually read from it. This matters enormously in this codebase, because it is what lets
\code{dashboard\_view} stack seven filters onto one QuerySet and still fire only one
query at the end.
\end{jargon}

\begin{jargon}{Migration}
A recorded, versioned instruction for changing the database's structure (adding a table,
adding a column). Django generates these from your models and applies them in order, so
that a fresh database and a long-lived one end up with identical structure.
\end{jargon}

\subsection{The directory layout}

\begin{figure}[H]
\centering
\begin{forest} dirtree
[paws-rescue-donations
  [manage.py {\rmfamily\itshape the command-line entry point}]
  [Procfile {\rmfamily\itshape tells the host how to start the app}]
  [requirements.txt {\rmfamily\itshape the six Python dependencies}]
  [.env.example {\rmfamily\itshape template for local configuration}]
  [db.sqlite3 {\rmfamily\itshape a committed database, see §\ref{sec:honest}}]
  [paws\_rescue/ {\rmfamily\itshape project configuration, not features}
    [settings.py {\rmfamily\itshape every configuration decision}]
    [urls.py {\rmfamily\itshape the top-level URL map}]
    [wsgi.py {\rmfamily\itshape the handle a web server grabs}]
  ]
  [apps/ {\rmfamily\itshape the two feature apps}
    [accounts/ {\rmfamily\itshape who may log in}
      [models.py] [forms.py] [views.py] [urls.py]
      [management/commands/create\_initial\_admin.py]
      [migrations/0001\_initial.py]
      [templates/accounts/]
      [tests.py]
    ]
    [donations/ {\rmfamily\itshape the actual subject matter}
      [models.py] [forms.py] [views.py] [urls.py]
      [migrations/0001\_initial.py]
      [templates/donations/]
      [tests.py]
    ]
  ]
  [templates/base.html {\rmfamily\itshape the page frame every page extends}]
  [static/css/custom.css]
]
\end{forest}
\caption{The tracked files, with the noise (empty \file{\_\_init\_\_.py} files,
screenshots) omitted.}
\end{figure}

\begin{jargon}{Django ``app''}
A confusing name. An app is not a separate program; it is a folder grouping the models,
views and templates for one area of responsibility. One Django \emph{project} contains
many apps. Here the project is \code{paws\_rescue} and it contains two apps,
\code{accounts} and \code{donations}.
\end{jargon}

\section{Architecture}

\subsection{The whole system at a glance}

\begin{figure}[H]
\centering
\resizebox{\textwidth}{!}{
\begin{tikzpicture}[node distance=8mm]
  \node[extbox] (browser) {Browser\\ \scriptsize Bootstrap 5 from a CDN};
  \node[box, below=12mm of browser] (mw) {Middleware chain\\ \scriptsize security, whitenoise, session, CSRF, auth};
  \node[box, below=10mm of mw] (urls) {\code{paws\_rescue/urls.py}\\ \scriptsize match the path};
  \node[box, font=\scriptsize, below left=12mm and 14mm of urls] (acc) {{\small\code{apps.accounts}}\\ login, logout,\\ admin management};
  \node[box, font=\scriptsize, below right=12mm and 14mm of urls] (don) {{\small\code{apps.donations}}\\ dashboard, CRUD,\\ CSV export};
  \node[box, below=26mm of urls] (orm) {Django ORM};
  \node[store, below=10mm of orm] (db) {PostgreSQL\\ or SQLite};
  \node[extbox, right=16mm of mw] (static) {Whitenoise\\ \scriptsize serves \file{custom.css}};

  \draw[flow] (browser) -- node[lbl,right]{HTTP request} (mw);
  \draw[flow] (mw) -- (urls);
  \draw[flow] (urls) -- (acc);
  \draw[flow] (urls) -- (don);
  \draw[flow] (acc) -- (orm);
  \draw[flow] (don) -- (orm);
  \draw[flow] (orm) -- node[lbl,right]{SQL} (db);
  \draw[dflow] (mw) -- (static);
  \draw[dflow] (mw.west) to[bend right=55] node[lbl,left]{rendered HTML} (browser.west);
\end{tikzpicture}}
\caption{Request path. Everything is server-rendered: the browser receives finished HTML
and runs no application JavaScript of the project's own.}
\end{figure}

\begin{jargon}{Middleware}
A stack of small components each request passes through on the way in, and each response
passes back through on the way out. They handle cross-cutting concerns so that no
individual view has to: attaching the logged-in user to the request, rejecting forged
form submissions, compressing static files.
\end{jargon}

\begin{jargon}{CDN (Content Delivery Network)}
A network of servers that hosts common files, here the Bootstrap CSS framework. The
page's HTML links out to \code{cdn.jsdelivr.net} rather than shipping Bootstrap in the
repository. It keeps the repository small; it also means the styling silently breaks
without an internet connection.
\end{jargon}

\subsection{The two apps and what separates them}

\begin{figure}[H]
\centering
\resizebox{0.95\textwidth}{!}{
\begin{tikzpicture}[node distance=6mm]
  \node[box, font=\scriptsize] (u) {{\small\textbf{User}}\\ UUID id, email (unique),\\ is\_staff, is\_superuser,\\ is\_active, created\_at};
  \node[box, font=\scriptsize, right=42mm of u] (d) {{\small\textbf{Donation}}\\ BigAuto id, donor\_name,\\ donor\_phone, donor\_email,\\ pet\_name, amount, currency,\\ payment\_method, reference\_no,\\ donation\_date, notes};
  \node[lbl, below=3mm of u] {table \code{users}};
  \node[lbl, below=3mm of d] {table \code{donations}};
  \draw[dflow] (u) -- node[lbl,above]{no foreign key exists\\ between these two} (d);
\end{tikzpicture}}
\caption{The complete data model. The dashed link is the relationship that is
\emph{absent}: nothing records which staff member entered a donation.}
\label{fig:model}
\end{figure}

\begin{honest}{The missing relationship is the design's biggest structural gap}
A donation-tracking system for a charity has no audit trail. \code{Donation} carries no
foreign key to \code{User}, so the database cannot answer ``who entered this record'' or
``who deleted it''. Deletes are hard deletes: \code{donation.delete()} removes the row
outright, with no tombstone and no history table. For a system whose entire purpose is
tracking other people's money, that is the first thing a reviewer will ask about. The
project's own README raises this too, under ``What I'd do differently''; it is recorded
here as confirmed against the schema in \file{apps/donations/migrations/0001\_initial.py},
which contains no \code{ForeignKey} operation.
\end{honest}

\begin{jargon}{Foreign key}
A column in one table holding the identifier of a row in another table, which is how
relational databases express ``this donation was entered by that user''. The database
then enforces that the referenced row actually exists.
\end{jargon}

\section{The \texttt{accounts} app: who is allowed in}

\subsection{Why there is a custom user model at all}

Django ships with a perfectly good \code{User} model. It has a \code{username} field, and
that is the problem: this application wants people to log in with an email address and
nothing else. Rather than carry a vestigial username column forever, the project defines
its own user model.

\begin{lstlisting}[language=Python,caption={\file{apps/accounts/models.py}, the model itself}]
class User(AbstractBaseUser, PermissionsMixin):
    id = models.UUIDField(primary_key=True, default=uuid.uuid4, editable=False)
    email = models.EmailField(unique=True, max_length=255)
    is_staff = models.BooleanField(default=False)
    is_superuser = models.BooleanField(default=False)
    is_active = models.BooleanField(default=True)
    created_at = models.DateTimeField(default=timezone.now)
    updated_at = models.DateTimeField(auto_now=True)

    objects = UserManager()

    USERNAME_FIELD = 'email'
    REQUIRED_FIELDS = []
\end{lstlisting}

Three decisions are packed into those lines.

\textbf{\code{AbstractBaseUser} plus \code{PermissionsMixin}.} \code{AbstractBaseUser}
gives password hashing and session handling but no fields beyond \code{password} and
\code{last\_login}, leaving the identity fields to you. \code{PermissionsMixin} adds the
\code{is\_superuser} flag and the group and permission relationships Django's own admin
site expects.

\textbf{\code{USERNAME\_FIELD = 'email'}} tells Django which field is the login
identifier. Everything downstream (the \code{authenticate()} function, the admin site)
reads this attribute rather than hard-coding ``username''.

\textbf{A UUID primary key.}

\begin{jargon}{Primary key, and UUID}
The primary key is the column uniquely identifying each row. The default is a
\emph{sequential integer}: 1, 2, 3. A UUID (Universally Unique Identifier) is instead a
random 128-bit value such as \code{09c139f5-1275-4a83-92c0-5684f5a58e8b}. Sequential
integers leak information (user \#3 can tell there are few users, and can guess that user
\#4 exists); UUIDs do not, and can be generated without asking the database.
\end{jargon}

\begin{honest}{The two primary key styles are inconsistent}
\code{User} uses a UUID. \code{Donation} uses Django's default sequential
\code{BigAutoField}, so donation URLs read \code{/donations/7/edit/}. Neither choice is
wrong on its own, but the project applies opposite reasoning to the two tables without
recording why. If the argument for the UUID was non-enumerability, it applies at least as
strongly to donation records, which are the sensitive data here. The honest reading,
labelled \textbf{inferred} because no comment in the source states a rationale, is that
the UUID came along with the hand-written user model and the default came along with the
generated one.
\end{honest}

\subsection{\texttt{UserManager}: why dropping \texttt{username} forces this}

Django's default user manager has a \code{create\_user(username, email, password)}
signature. With no username field, that signature no longer fits, and every part of
Django that calls it (notably \code{manage.py createsuperuser}) would break. The project
supplies a replacement.

\begin{lstlisting}[language=Python,caption={\file{apps/accounts/models.py}, the manager}]
class UserManager(BaseUserManager):
    def create_user(self, email, password=None, **extra_fields):
        if not email:
            raise ValueError('The Email field must be set')
        email = self.normalize_email(email)
        user = self.model(email=email, **extra_fields)
        user.set_password(password)
        user.save(using=self._db)
        return user

    def create_superuser(self, email, password=None, **extra_fields):
        extra_fields.setdefault('is_staff', True)
        extra_fields.setdefault('is_superuser', True)
        if extra_fields.get('is_staff') is not True:
            raise ValueError('Superuser must have is_staff=True.')
        if extra_fields.get('is_superuser') is not True:
            raise ValueError('Superuser must have is_superuser=True.')
        return self.create_user(email, password, **extra_fields)
\end{lstlisting}

Two details carry real weight. \code{normalize\_email} lower-cases the domain part of
the address, so \code{Admin@EXAMPLE.com} and \code{Admin@example.com} become the same
string and the \code{unique=True} constraint actually bites. And \code{set\_password}
hashes the password rather than storing it:

\begin{jargon}{Password hashing}
Storing a one-way scrambled form of a password instead of the password itself. Django's
default here is PBKDF2 with SHA-256 and 600{,}000 iterations. Given the stored value you
cannot practically recover the password, and the deliberate slowness makes guessing
expensive. \code{set\_password} does this; assigning to \code{user.password} directly
would store plain text, which is exactly the mistake the manager exists to prevent.
\end{jargon}

\begin{honest}{\texttt{create\_user} accepts \texttt{password=None}}
Calling \code{create\_user(email='x@y.com')} with no password is legal. Django's
\code{set\_password(None)} produces an unusable password rather than an empty one, so
this does not create a login bypass, and the account simply cannot be logged into. It is
worth knowing rather than fixing, because a reader will notice the default and wonder.
\end{honest}

\subsection{The login path}

\begin{figure}[H]
\centering
\resizebox{\textwidth}{!}{
\begin{tikzpicture}[node distance=7mm]
  \node[box] (post) {POST to \code{/login/}\\ \scriptsize email + password};
  \node[dec, below=8mm of post] (auth) {\code{authenticate()}\\ succeeds?};
  \node[box, left=16mm of auth] (bad) {add error message,\\ re-render form};
  \node[dec, below=10mm of auth] (staff) {\code{user.is\_staff}?};
  \node[box, left=16mm of staff] (nostaff) {``You do not have\\ permission''};
  \node[box, below=10mm of staff] (ok) {\code{login(request, user)}\\ \scriptsize session cookie set};
  \node[box, below=8mm of ok] (dash) {redirect to \code{/dashboard/}};

  \draw[flow] (post) -- (auth);
  \draw[flow] (auth) -- node[lbl,above]{no} (bad);
  \draw[flow] (auth) -- node[lbl,right]{yes} (staff);
  \draw[flow] (staff) -- node[lbl,above]{no} (nostaff);
  \draw[flow] (staff) -- node[lbl,right]{yes} (ok);
  \draw[flow] (ok) -- (dash);
\end{tikzpicture}}
\caption{\code{login\_view}. The \code{is\_staff} gate is checked \emph{after}
authentication succeeds, and a non-staff user is never actually logged in.}
\end{figure}

\code{EmailAuthenticationForm} does the authentication inside its \code{clean()} method
and stashes the result on \code{self.user\_cache}, which \code{get\_user()} returns. This
mirrors how Django's own \code{AuthenticationForm} works. Note the small awkwardness the
custom user model forces:

\begin{lstlisting}[language=Python,caption={\file{apps/accounts/forms.py}, inside \code{clean}}]
self.user_cache = authenticate(
    self.request,
    username=email,      # the keyword is still 'username'
    password=password
)
\end{lstlisting}

\begin{keyidea}{Why \texttt{username=email} is correct, not a bug}
\code{authenticate()} passes its keyword arguments through to the authentication backend,
and Django's \code{ModelBackend} looks up whatever field \code{USERNAME\_FIELD} names
using the value it received under the key \code{username}. The keyword is a fixed part of
the interface; the value is an email. The same wrinkle shows up in the tests, which call
\code{self.client.login(username='admin@test.com', ...)}.
\end{keyidea}

\subsection{Admin management, and where it is weaker than it looks}

\code{admin\_list\_view} lists every \code{is\_staff} user. \code{admin\_create\_view}
creates one. \code{admin\_delete\_view} deletes one, refusing to delete the requester.

\begin{lstlisting}[language=Python,caption={\file{apps/accounts/views.py}, self-deletion guard}]
user_to_delete = get_object_or_404(User, id=user_id)

if user_to_delete.id == request.user.id:
    messages.error(request, 'You cannot delete your own account.')
    return redirect('admin_list')
\end{lstlisting}

The guard is enforced in the view, not merely hidden in the template, which is the right
way round: hiding a button does not stop anyone who can send an HTTP request. The
template hides the button too, purely as a courtesy. This is covered by
\code{test\_cannot\_delete\_self}.

\begin{honest}{Every ``admin'' created through the web UI is a full superuser}
\code{AdminCreationForm.save()} hard-codes both flags:
\begin{lstlisting}[language=Python]
user.is_staff = True
user.is_superuser = True
\end{lstlisting}
There is no way to create a limited account through the interface, so the ``admin
management'' feature has exactly one privilege level. The \code{PermissionsMixin}, groups
and per-object permissions are all wired up and completely unused. For a rescue where a
volunteer might reasonably be allowed to add donations but not to delete other logins,
the distinction the model already supports is thrown away at the form layer.
\end{honest}

\begin{honest}{The password validators do not apply to admin creation, verified by running it}
\file{settings.py} configures the four standard \code{AUTH\_PASSWORD\_VALIDATORS}
(similarity, minimum length, common-password, numeric-only). Those validators only run
where something explicitly calls them, and \code{AdminCreationForm} never does: it
subclasses \code{forms.ModelForm}, not Django's \code{UserCreationForm}. Confirmed by
executing the form directly against the project's settings:
\begin{lstlisting}
AdminCreationForm({'email': 'weak@example.com', 'password': '1'})
  -> is_valid() == True, errors == {}
\end{lstlisting}
A one-character password is accepted for an account that is, per the finding above, a
full superuser. The validators do run for \code{manage.py createsuperuser}, which is
precisely why the gap is easy to miss: the configuration looks present and correct at the
top of \file{settings.py}. The fix is one line, calling
\code{django.contrib.auth.password\_validation.validate\_password} inside a
\code{clean\_password} method.
\end{honest}

\begin{honest}{Logging out is possible via GET}
\code{logout\_view} is decorated \code{@require\_http\_methods(["POST", "GET"])}. The
navigation bar submits a proper POST form with a CSRF token, so normal use is fine, but
the GET route remains open, which means any page anywhere can log the user out by
embedding \code{<img src="https://thesite/logout/">}. This is a nuisance rather than a
breach (no data is exposed and nothing is destroyed), and Django removed GET logout from
its own \code{LogoutView} in version 4.1 for exactly this reason. Dropping \code{"GET"}
from the list costs nothing, since the only caller already uses POST.
\end{honest}

\begin{jargon}{CSRF (Cross-Site Request Forgery)}
An attack where a page you are visiting quietly submits a request to a different site you
happen to be logged into, riding on your session cookie. Django's defence is a secret
token embedded in every form (\code{\{\% csrf\_token \%\}}) and checked on submission;
the attacking page cannot read the token, so its forged request fails. It only protects
methods that change state, which is why leaving a state-changing action reachable by GET
sidesteps the whole mechanism.
\end{jargon}

\subsection{The \texttt{create\_initial\_admin} command}

A chicken-and-egg problem: the app has no public sign-up, so a fresh deployment has no
users and nobody can log in to create one. Django's interactive
\code{manage.py createsuperuser} cannot help on a hosting platform with no terminal
attached. The project adds a custom management command reading two environment
variables.

\begin{lstlisting}[language=Python,caption={\file{apps/accounts/management/commands/create\_initial\_admin.py}, the core}]
email = os.getenv('INITIAL_ADMIN_EMAIL')
password = os.getenv('INITIAL_ADMIN_PASSWORD')

if not email or not password:
    self.stdout.write(self.style.WARNING('... must be set ...'))
    return

if User.objects.filter(email=email).exists():
    self.stdout.write(self.style.WARNING(f'... already exists. Skipping.'))
    return

User.objects.create_superuser(email=email, password=password)
\end{lstlisting}

The existence check makes the command \emph{idempotent}, which is what makes it safe to
wire into a deploy script that runs on every release.

\begin{jargon}{Idempotent}
Running it once and running it ten times produce the same result. Here, the second run
finds the user already present and does nothing, rather than crashing on the unique-email
constraint or resetting the password.
\end{jargon}

\begin{honest}{Both failure modes exit successfully and quietly}
Missing variables print a warning and \code{return}. The command still exits with status
code 0, meaning ``success''. A deploy pipeline running
\code{python manage.py create\_initial\_admin} and checking the exit code sees a green
tick and moves on, and the operator discovers there is no admin account only when nobody
can log in. Raising \code{CommandError} instead would make the deploy fail loudly, which
is the behaviour you want.
\end{honest}

\section{The \texttt{donations} app: the actual subject matter}

\subsection{The \texttt{Donation} model}

\begin{lstlisting}[language=Python,caption={\file{apps/donations/models.py}, abridged}]
class Donation(models.Model):
    PAYMENT_METHOD_CHOICES = [
        ('cash', 'Cash'), ('bank', 'Bank Transfer'),
        ('online', 'Online Payment'), ('other', 'Other'),
    ]

    donor_name  = models.CharField(max_length=255)
    donor_phone = models.CharField(max_length=50, blank=True, null=True)
    donor_email = models.EmailField(blank=True, null=True)
    pet_name    = models.CharField(max_length=255, help_text='Name of the pet ...')
    amount      = models.DecimalField(max_digits=10, decimal_places=2,
                      validators=[MinValueValidator(Decimal('0.00'))])
    currency    = models.CharField(max_length=10, default='PKR')
    payment_method = models.CharField(max_length=20,
                      choices=PAYMENT_METHOD_CHOICES, default='cash')
    reference_no   = models.CharField(max_length=255, blank=True, null=True)
    donation_date  = models.DateTimeField(default=timezone.now)
    notes          = models.TextField(blank=True, null=True)
    created_at     = models.DateTimeField(auto_now_add=True)
    updated_at     = models.DateTimeField(auto_now=True)

    class Meta:
        db_table = 'donations'
        ordering = ['-donation_date']
\end{lstlisting}

\begin{keyidea}{\texttt{DecimalField}, and why money is never a float}
\code{amount} is a \code{DecimalField}, not a \code{FloatField}. Floating-point numbers
store values in binary, and 0.1 has no exact binary representation, so \code{0.1 + 0.2}
famously yields \code{0.30000000000000004}. Summing thousands of donations that way
accumulates real error. A decimal stores the digits exactly.
\code{max\_digits=10, decimal\_places=2} means up to eight digits before the point and
exactly two after: the largest recordable donation is 99{,}999{,}999.99.
\end{keyidea}

\begin{keyidea}{The difference between \texttt{blank} and \texttt{null}}
They look redundant and are not. \code{null=True} is a \emph{database} statement: this
column may hold SQL NULL. \code{blank=True} is a \emph{validation} statement: forms may
leave this field empty. Django's convention is to avoid \code{null=True} on text fields,
because it creates two different ways to say ``no value'' (an empty string and NULL) that
must then be handled separately forever. This project sets both on every optional text
field, so \code{donor\_phone} can be either \code{''} or \code{None}, which is why the
CSV writer has to say \code{donation.donor\_phone or ''} on every such column.
\end{keyidea}

\code{auto\_now\_add} stamps \code{created\_at} once at insert; \code{auto\_now} refreshes
\code{updated\_at} on every save. \code{donation\_date} is separate from both because
when the money arrived and when somebody typed it in are different facts.

\code{Meta.ordering = ['-donation\_date']} makes newest-first the default for every query
that does not say otherwise. \code{db\_table = 'donations'} overrides Django's generated
name (\code{donations\_donation}), which is a cosmetic preference.

\begin{honest}{\texttt{currency} is free text, and the consequence is worse than the README says}
\code{currency} is a \code{CharField(max\_length=10, default='PKR')} with no
\code{choices}. Anything up to ten characters is accepted. The README's own
self-assessment says typos such as ``pkr'' or ``Rs'' would ``silently split the totals''.
Running the code shows the opposite, and the opposite is more dangerous. See the
dedicated finding in §\ref{sec:currencybug}.
\end{honest}

\subsection{Two forms, doing quite different jobs}

\code{DonationForm} is a \code{ModelForm}: it derives its fields from the model, so the
field list is declared once. Its only additions are Bootstrap CSS classes on each widget
and two \code{clean} methods.

\begin{lstlisting}[language=Python,caption={\file{apps/donations/forms.py}, the two cleaners}]
def clean_amount(self):
    amount = self.cleaned_data.get('amount')
    if amount and amount < 0:
        raise forms.ValidationError('Amount must be greater than or equal to 0.')
    return amount

def clean_donation_date(self):
    """Set donation_date to now if it's empty."""
    donation_date = self.cleaned_data.get('donation_date')
    if not donation_date:
        from django.utils import timezone
        donation_date = timezone.now()
    return donation_date
\end{lstlisting}

\code{clean\_donation\_date} implements ``leave the field empty to mean now''. The model
default (\code{default=timezone.now}) only applies when the field is omitted entirely;
a form submitting an empty string is not an omission, so without this method an empty
date box would be a validation error rather than a convenience.

\begin{honest}{\texttt{clean\_amount} is a redundant second guard, and I checked which one actually fires}
There are two negative-amount defences: \code{MinValueValidator} on the model field and
\code{clean\_amount} on the form. I instrumented the form and ran both paths against the
project's real settings:
\begin{itemize}
  \item With \code{clean\_amount} present, submitting \code{-100.00} is rejected with
    \emph{its} message: ``Amount must be greater than or equal to 0.''
  \item With \code{clean\_amount} neutralised, \code{-100.00} is \emph{still} rejected,
    with the model validator's message: ``Ensure this value is greater than or equal to
    0.00.''
\end{itemize}
So the method is not dead code, but it is not load-bearing either: it only changes the
wording. The reason both fire at all is worth knowing, because it is a genuine Django
subtlety. Model-field validators are \textbf{not} copied onto the form field (I printed
them: the form field carries only a \code{DecimalValidator}). They run later, in
\code{ModelForm.\_post\_clean()}, which calls \code{instance.full\_clean()}. That is
after every \code{clean\_<field>} method. Hence \code{clean\_amount} gets first refusal
and its message wins.
\end{honest}

\begin{honest}{The \texttt{if amount and amount < 0} condition is subtly wrong, harmlessly}
\code{Decimal('0.00')} is falsy in Python, so a zero amount short-circuits the
\code{and} and skips the comparison entirely. It happens not to matter, because zero is
not less than zero and should pass anyway. Verified: a \code{0.00} donation is accepted
by the form. The condition should be \code{if amount is not None and amount < 0}; the
same \code{is not None} care is taken correctly a few files away, in the amount filters
in \code{dashboard\_view}. The codebase knows the idiom and applies it inconsistently.
\end{honest}

\code{DonationFilterForm} is a plain \code{forms.Form}, not tied to the model. Every one
of its eight fields is \code{required=False}, which is what makes ``filter by any
combination of these'' work. It exists to do type coercion and nothing else: it turns the
raw query string \code{?amount\_min=500\&date\_from=2026-01-01} into a validated
\code{Decimal} and a real \code{date} object, or into an error.

\subsection{\texttt{dashboard\_view}: the heart of the application}

This one function does filtering, sorting, aggregation, pagination and session-writing.
It is the most consequential code in the project.

\begin{figure}[H]
\centering
\resizebox{\textwidth}{!}{
\begin{tikzpicture}[node distance=6mm]
  \node[box] (all) {\code{Donation.objects.all()}\\ \scriptsize lazy: no query yet};
  \node[box, below=7mm of all] (f1) {\code{search} $\rightarrow$ Q(donor\_name) OR Q(pet\_name)};
  \node[box, below=5mm of f1] (f2) {\code{date\_from} / \code{date\_to} $\rightarrow$ \code{donation\_date\_\_date} range};
  \node[box, below=5mm of f2] (f3) {\code{payment\_method} $\rightarrow$ exact};
  \node[box, below=5mm of f3] (f4) {\code{currency} $\rightarrow$ \code{iexact}};
  \node[box, below=5mm of f4] (f5) {\code{amount\_min} / \code{amount\_max} $\rightarrow$ gte / lte};
  \node[box, below=5mm of f5] (sort) {\code{order\_by(sort\_by)}};

  \node[box, right=30mm of f1] (agg1) {\code{total\_all\_time}\\ \scriptsize from \code{objects.all()}, unfiltered};
  \node[box, below=6mm of agg1] (agg2) {\code{total\_filtered}\\ \scriptsize from the filtered set};
  \node[box, below=6mm of agg2] (agg3) {\code{top\_pets}\\ \scriptsize values+annotate on filtered set};
  \node[box, below=6mm of agg3] (page) {\code{Paginator(donations, 20)}};
  \node[store, below=6mm of page] (sess) {session\\ \scriptsize \code{filtered\_donation\_ids}};

  \draw[flow] (all) -- (f1); \draw[flow] (f1) -- (f2); \draw[flow] (f2) -- (f3);
  \draw[flow] (f3) -- (f4); \draw[flow] (f4) -- (f5); \draw[flow] (f5) -- (sort);
  \draw[flow] (sort.east) to[bend right=20] node[lbl,below]{one QuerySet,\\ now evaluated} (agg2.west);
  \draw[flow] (agg2) -- (agg3); \draw[flow] (agg3) -- (page); \draw[flow] (page) -- (sess);
\end{tikzpicture}}
\caption{\code{dashboard\_view}. Each filter is applied only if its cleaned value is
present, so the seven are freely combinable. Note \code{total\_all\_time} deliberately
branches off the \emph{unfiltered} manager.}
\label{fig:dash}
\end{figure}

\begin{keyidea}{Why stacking filters like this is cheap}
Each \code{donations = donations.filter(...)} line looks like it runs a query. None of
them do. A QuerySet is lazy: \code{filter()} returns a new QuerySet with one more
condition recorded, and the database is not touched until something iterates the results.
All seven conditions collapse into a single \code{WHERE} clause in one SQL statement.
\end{keyidea}

\begin{jargon}{\texttt{Q} object}
Django's way of expressing a condition that is not a simple AND. Keyword arguments to
\code{filter()} are ANDed together, so there is no way to spell OR with them. A \code{Q}
object wraps a condition so it can be combined with \code{|} (OR) and \code{\&} (AND).
The search box needs \code{Q(donor\_name\_\_icontains=s) | Q(pet\_name\_\_icontains=s)}:
match either column.
\end{jargon}

\begin{jargon}{Field lookups: \texttt{icontains}, \texttt{iexact}, \texttt{gte}, \texttt{lte}}
The double-underscore suffix on a filter keyword chooses the comparison.
\code{icontains} is a case-insensitive substring match (SQL \code{LIKE '\%x\%'}),
\code{iexact} a case-insensitive equality, \code{gte} and \code{lte} are ``greater than
or equal'' and ``less than or equal''. \code{donation\_date\_\_date\_\_gte} chains two:
extract the date part of the timestamp, then compare.
\end{jargon}

\begin{keyidea}{Why \texttt{donation\_date\_\_date} matters}
\code{donation\_date} is a timestamp, not a date. Filtering \code{donation\_date\_\_gte}
against a date silently means midnight on that date, so a donation at 14:00 on the
\code{date\_to} day would be excluded from a range that visibly includes that day. The
\code{\_\_date} lookup extracts the calendar date first and compares dates to dates, which
is what a user picking ``to: 7 February'' means.
\end{keyidea}

The amount filters show the careful version of the idiom the form got wrong:

\begin{lstlisting}[language=Python,caption={\file{apps/donations/views.py}, the amount bounds}]
amount_min = filter_form.cleaned_data.get('amount_min')
if amount_min is not None:
    donations = donations.filter(amount__gte=amount_min)
\end{lstlisting}

\code{is not None} rather than a truthiness test, because \code{Decimal('0')} is falsy
and ``show me donations of at least 0'' is a legitimate request that a bare
\code{if amount\_min:} would silently discard.

\subsubsection{The aggregation, and a trap it avoids by luck}
\label{sec:groupby}

\begin{lstlisting}[language=Python,caption={\file{apps/donations/views.py}, analytics}]
total_all_time = Donation.objects.aggregate(total=Sum('amount'))['total'] or 0
total_filtered = donations.aggregate(total=Sum('amount'))['total'] or 0

top_pets = (
    donations
    .values('pet_name')
    .annotate(total_amount=Sum('amount'), donation_count=Count('id'))
    .order_by('-total_amount')[:5]
)
\end{lstlisting}

\begin{jargon}{\texttt{values} plus \texttt{annotate}, i.e. GROUP BY}
\code{.values('pet\_name')} says ``I care only about this column''. Following it with
\code{.annotate(Sum('amount'))} makes Django group the rows by that column and compute
the sum within each group. It is SQL's \code{GROUP BY pet\_name} with
\code{SUM(amount)}, spelled in Python.
\end{jargon}

The README records that this panel once aggregated over \code{Donation.objects.all()} and
therefore contradicted the filtered total displayed beside it. The current code
aggregates over the filtered \code{donations} QuerySet, which is correct, and I confirmed
the fix holds.

There is a well-known Django trap lurking here that the code does \emph{not} fall into,
and it is worth recording as verified rather than assumed. When a QuerySet carries an
\code{order\_by}, and this one always does (\code{Meta.ordering}, plus an explicit
\code{order\_by(sort\_by)} a few lines above), the ordering columns can be dragged into
the \code{GROUP BY} clause, shattering each pet into one group per distinct timestamp. I
tested it directly: two donations for the pet Max on different dates, totalling 400.

\begin{lstlisting}[caption={Observed output, all three sort options}]
top_pets (default sort):        [{'pet_name': 'Max', 'total_amount': 400, 'donation_count': 2}, ...]
top_pets (sort_by=-donation_date): [{'pet_name': 'Max', 'total_amount': 400, 'donation_count': 2}, ...]
top_pets (sort_by=-amount):     [{'pet_name': 'Max', 'total_amount': 400, 'donation_count': 2}, ...]

generated SQL:
SELECT "donations"."pet_name", CAST(SUM("donations"."amount") AS NUMERIC) AS "total_amount",
       COUNT("donations"."id") AS "donation_count"
FROM "donations" GROUP BY "donations"."pet_name" ORDER BY 2 DESC LIMIT 5
\end{lstlisting}

One row for Max, 400, and a \code{GROUP BY} on \code{pet\_name} alone. The trailing
sort on \code{-total\_amount} is what saves it: an \code{order\_by} \emph{replaces} any
previous ordering rather than adding to it, so the inherited date ordering is discarded
before the SQL is built. The aggregation is correct. It is correct because a clause
added for a different reason, sorting the panel by size, happens to clear the ordering
that would have broken it. That is worth knowing if anyone ever removes that line.

\begin{keyidea}{The \texttt{or 0} is load-bearing}
\code{Sum} over zero rows returns \code{None}, not \code{0}. Without \code{or 0} an empty
database would render ``None PKR'' on the dashboard.
\end{keyidea}

\begin{jargon}{Pagination}
Splitting a long list across numbered pages. \code{Paginator(donations, 20)} slices the
QuerySet into 20-row pages and, because the QuerySet is lazy, does it with SQL
\code{LIMIT} and \code{OFFSET} rather than loading every donation into memory.
\code{get\_page} is the forgiving variant: a missing, non-numeric or out-of-range page
number yields the first or last page instead of raising.
\end{jargon}

\subsection{CSV export, and the two bugs in it}

\begin{lstlisting}[language=Python,caption={\file{apps/donations/views.py}, the session handoff}]
# in dashboard_view, on every render:
request.session['filtered_donation_ids'] = list(donations.values_list('id', flat=True))

# in export_csv_view:
filtered_ids = request.session.get('filtered_donation_ids', [])

if filtered_ids:
    donations = Donation.objects.filter(id__in=filtered_ids).order_by('-donation_date')
else:
    donations = Donation.objects.all().order_by('-donation_date')
\end{lstlisting}

The design problem it is solving is real. ``Export CSV'' is a plain link: a GET to the
export URL, carrying no filter state at all, so the export view has no idea what the
dashboard was showing. The chosen workaround: the dashboard writes the id of every
matching donation into the session, and the export view reads them back.

\begin{jargon}{Session}
Server-side storage tied to one browser via a cookie. The browser holds only an opaque
key; the data (here, a list of donation ids) lives in the server's \code{django\_session}
database table.
\end{jargon}

The README calls this ``the weakest part of the design''. That is correct, and running it
surfaces two concrete defects that the README does not name.

\begin{honest}{A filter matching nothing exports the entire database, verified by running it}
\code{if filtered\_ids:} cannot distinguish ``the user has not loaded the dashboard, so
export everything'' from ``the user's filter matched zero rows, so export nothing''. Both
are the empty list, and the empty list is falsy, so both take the \code{else} branch. I
drove this through the real test client with two donations in the database:
\begin{lstlisting}
GET /dashboard/?search=NoSuchDonorZZZ
  -> dashboard displays "0 donation(s)"
  -> session['filtered_donation_ids'] == []
GET /donations/export/
  -> CSV data rows exported: 2
  -> contains John? True | contains Jane? True
\end{lstlisting}
The user filters to nothing, sees an empty table reading ``No donations found'', clicks
Export CSV, and silently downloads every donor name, phone number and email address in
the system. This is the most serious defect in the project: it is a real bug, it is
triggered by an ordinary action, it fails in the unsafe direction, and it is invisible
because the file downloads without complaint. The fix is one character of intent:
\code{if filtered\_ids is not None:} combined with storing \code{None} rather than
never writing the key, or better, re-deriving the queryset from the query string.
\end{honest}

\begin{honest}{The export is stale by construction}
The session holds ids captured at the moment the dashboard was last rendered. Anything
that happens afterwards is invisible to it. Verified: filter to donor ``John'' (session
holds \code{[1]}), delete John's donation in another window, then export. The
\code{id\_\_in} lookup simply finds nothing and the CSV comes back with 0 data rows, so
this particular case degrades safely. A donation \emph{added} after the render is
likewise missing from an export the user believes reflects their filter. There is also no
bound on growth: every dashboard render writes one integer per matching row into the
session, and with \code{SESSION\_SAVE\_EVERY\_REQUEST = True} in \file{settings.py} that
row is rewritten to the database on every single request. At a few hundred donations this
is invisible; it is still the wrong shape.
\end{honest}

\begin{keyidea}{The fix the README already identifies}
Re-apply \code{DonationFilterForm} to \code{request.GET} inside \code{export\_csv\_view}
and rebuild the queryset from the query string, with the export link carrying the current
filters. No session, no staleness, no growth, no empty-list ambiguity. It needs the
filtering logic extracted into a shared helper, which the README also flags: right now
\code{dashboard\_view} applies the seven filters inline and nothing else can reuse them.
\end{keyidea}

The writing itself is unremarkable and correct: \code{HttpResponse(content\_type='text/csv')},
a \code{Content-Disposition} header naming the file
\code{donations\_YYYY-MM-DD.csv} to make the browser download rather than display it, and
\code{csv.writer(..., quoting=csv.QUOTE\_ALL)}.

\begin{keyidea}{Why \texttt{QUOTE\_ALL} is the right call}
CSV separates fields with commas, so a value containing a comma (``Bella, the tabby'', or
any note with a comma in it) would otherwise split into two columns and corrupt every
column after it. Quoting every field unconditionally makes that impossible.
\end{keyidea}

\code{get\_payment\_method\_display()} is worth pointing out: for any field with
\code{choices}, Django generates a \code{get\_<field>\_display()} method returning the
human label (``Bank Transfer'') rather than the stored value (``bank''). The CSV uses it;
so does the dashboard template.

\subsection{The three CRUD views}

\begin{jargon}{CRUD}
Create, Read, Update, Delete: the four basic operations on a record.
\end{jargon}

\code{donation\_create\_view} and \code{donation\_update\_view} are the standard Django
form pattern, differing only by \code{instance=donation}: passing an existing object
makes the same \code{ModelForm} an edit form rather than a create form.
\code{donation\_delete\_view} shows a confirmation page on GET and deletes on POST, which
is the correct division (a GET must never change state).

\begin{honest}{Form errors are pushed through the messages framework, twice over}
Both write views do this on failure:
\begin{lstlisting}[language=Python]
for field, errors in form.errors.items():
    for error in errors:
        messages.error(request, f'{field}: {error}')
\end{lstlisting}
The bound form is then re-rendered, and the template also displays each field's errors
inline. So a validation failure shows the user the same message twice, once in a red
banner reading \code{donor\_name: This field is required.} with the raw field name, and
once beside the field. The banner version leaks the internal attribute name into the UI.
Deleting the loop loses nothing.
\end{honest}

\begin{honest}{\texttt{is\_admin} is defined twice}
The identical four-line function exists in both \file{apps/accounts/views.py} and
\file{apps/donations/views.py}. Harmless today, and exactly how two definitions drift
apart later.
\end{honest}

\section{Configuration: \texttt{settings.py}}

\subsection{Environment variables and the twelve-factor idea}

\begin{lstlisting}[language=Python,caption={\file{paws\_rescue/settings.py}, the top}]
load_dotenv()

SECRET_KEY = os.getenv('SECRET_KEY', 'django-insecure-dev-key-change-in-production')
DEBUG = os.getenv('DEBUG', 'False') == 'True'
ALLOWED_HOSTS = os.getenv('ALLOWED_HOSTS', 'localhost,127.0.0.1').split(',')
TESTING = 'test' in sys.argv
\end{lstlisting}

\begin{jargon}{Environment variable, and \texttt{.env}}
A named value handed to a program by the system that starts it, rather than written in
its source. Secrets live there so they never enter version control. \code{python-dotenv}
reads a local \file{.env} file into the environment for development convenience;
\file{.env} is in \file{.gitignore}, and \file{.env.example} documents the required keys
with placeholder values.
\end{jargon}

\code{DEBUG = os.getenv('DEBUG', 'False') == 'True'} is a string comparison because
environment variables are always strings: the literal text \code{True} is the only value
that enables debug mode, and \code{1}, \code{true} or \code{yes} all leave it off.
Defaulting to \code{False} is the right way round: a forgotten variable produces a secure
app, not an insecure one.

\begin{honest}{The \texttt{SECRET\_KEY} fallback silently produces a working, insecure app}
\code{SECRET\_KEY} falls back to a hard-coded string in the source. The key signs session
cookies and CSRF tokens; anyone who knows it can forge either. The word
\code{django-insecure} in the value is the only warning, it is committed to a public
repository, and a deployment that forgets to set \code{SECRET\_KEY} starts up perfectly
happily using it. \code{DEBUG} was given a safe default; \code{SECRET\_KEY} was not.
Crashing on a missing key in production is the standard remedy.
\end{honest}

\subsection{The database switch}

\begin{lstlisting}[language=Python,caption={\file{paws\_rescue/settings.py}, database selection}]
DATABASE_URL_ENV = os.getenv('DATABASE_URL')

if DATABASE_URL_ENV:
    DATABASES = {'default': dj_database_url.config(
        default=DATABASE_URL_ENV, conn_max_age=600, conn_health_checks=True)}
else:
    DATABASES = {'default': {
        'ENGINE': 'django.db.backends.sqlite3', 'NAME': BASE_DIR / 'db.sqlite3'}}
\end{lstlisting}

Set \code{DATABASE\_URL} and the app talks to PostgreSQL; leave it empty and it uses a
local SQLite file. This is what lets a new contributor run the project with no database
server installed.

\begin{jargon}{SQLite versus PostgreSQL}
SQLite is a database that is just a file on disk, with no server process. It is ideal for
development and hopeless for a multi-user production deployment, where concurrent writes
serialise behind a single lock and the file vanishes with the container's disk.
PostgreSQL is a full server process built for exactly that.
\end{jargon}

A \code{conn\_max\_age} of 600 keeps a database connection alive for ten minutes rather
than reconnecting on every request, since a TCP handshake plus authentication on each
page load is real latency. Setting \code{conn\_health\_checks} to \code{True} is the
necessary companion: it pings a reused connection before handing it over, so a
connection the database closed in the meantime produces a reconnect, not an error page.

\subsection{Static files and Whitenoise}

\begin{jargon}{Static files}
Files served exactly as they are, with no per-request computation: CSS, JavaScript,
images. Django serves them itself in development and refuses to in production, on the
grounds that a real web server does it better.
\end{jargon}

\begin{jargon}{Whitenoise}
A library letting the Python application serve its own static files competently, with
compression and cache headers, so a deployment needs no separate nginx. Its middleware
sits second in the chain, immediately after Django's security middleware, so a request
for a CSS file is answered and returned without waking the rest of the stack.
\end{jargon}

\begin{jargon}{\texttt{collectstatic} and manifest storage}
\code{collectstatic} copies every static file from every app into one directory
(\code{STATIC\_ROOT}) for serving. \code{CompressedManifestStaticFilesStorage} goes
further: it renames each file to include a hash of its contents
(\code{custom.a1b2c3.css}) and writes a manifest mapping original names to hashed ones.
This allows caching forever, since changed content means a changed filename. The cost is
that \code{\{\% static 'css/custom.css' \%\}} now performs a manifest lookup and
\textbf{raises an exception} if the entry is missing.
\end{jargon}

That cost is precisely what produced one of the project's documented battles:

\begin{lstlisting}[language=Python,caption={\file{paws\_rescue/settings.py}, the test carve-out}]
if TESTING:
    STORAGES = {"staticfiles": {
        "BACKEND": "django.contrib.staticfiles.storage.StaticFilesStorage"}}
else:
    STORAGES = {"staticfiles": {
        "BACKEND": "whitenoise.storage.CompressedManifestStaticFilesStorage"}}

# ... and at the bottom:
if not DEBUG and not TESTING:
    SECURE_SSL_REDIRECT = True
    SESSION_COOKIE_SECURE = True
    CSRF_COOKIE_SECURE = True
    SECURE_BROWSER_XSS_FILTER = True
    SECURE_CONTENT_TYPE_NOSNIFF = True
    X_FRAME_OPTIONS = 'DENY'
\end{lstlisting}

\begin{keyidea}{Why \texttt{TESTING = 'test' in sys.argv} has to exist}
This is the single most instructive line in the settings file. \code{manage.py test} does
\emph{not} set \code{DEBUG=True}; it runs with whatever \code{DEBUG} the environment
gives, which defaults to \code{False}. So a block guarded only by \code{if not DEBUG}
applies during tests, and two things then break at once:
\begin{itemize}
  \item \code{SECURE\_SSL\_REDIRECT} turns every test-client request into a 301 redirect
    to \code{https://}, so every assertion expecting a 200 fails.
  \item The manifest storage raises on the first \code{\{\% static \%\}} tag, because
    tests never run \code{collectstatic}.
\end{itemize}
Inspecting \code{sys.argv} for the string \code{test} is admittedly a hack, and it is the
common one. The clean alternative is a separate settings module for tests.
\end{keyidea}

\begin{jargon}{The production security flags, one at a time}
\code{SECURE\_SSL\_REDIRECT} bounces plain HTTP to HTTPS. \code{SESSION\_COOKIE\_SECURE}
and \code{CSRF\_COOKIE\_SECURE} tell the browser never to transmit those cookies over
plain HTTP, so a network eavesdropper cannot lift a session.
\code{SECURE\_CONTENT\_TYPE\_NOSNIFF} stops browsers from guessing a file's type and
executing an upload as script. \code{X\_FRAME\_OPTIONS = 'DENY'} forbids embedding the
site in an iframe, defeating clickjacking. \code{SECURE\_BROWSER\_XSS\_FILTER} sets a
legacy header that modern browsers ignore entirely; it is harmless but does nothing.
\end{jargon}

\begin{honest}{The security block being conditional is a real fragility}
These flags apply only when \code{DEBUG} is false and tests are not running. The
condition works, but it means the production security posture is determined by two
unrelated inferences (an environment string, and the contents of \code{sys.argv}) rather
than declared. The README names the fix, splitting settings into base, dev and prod
modules, and this is the honest assessment of why: today, anyone who runs the app with
\code{DEBUG=True} on a public host loses all six protections silently and simultaneously.
\end{honest}

\subsection{Sessions}

\begin{lstlisting}[language=Python]
SESSION_COOKIE_AGE = 86400          # 1 day
SESSION_SAVE_EVERY_REQUEST = True
\end{lstlisting}

\code{SESSION\_SAVE\_EVERY\_REQUEST} rewrites the session row on every request, which
makes the one-day expiry a \emph{sliding} window: an active user is never logged out
mid-task, and an idle one is logged out a day later. The cost is a database write on
every single request, which combines badly with the CSV export writing the whole id list
into that session on every dashboard render.

\section{Templates and the front end}

\file{templates/base.html} defines the frame: Bootstrap 5 and Bootstrap Icons from a CDN,
the project's \file{custom.css}, a navigation bar rendered only when
\code{user.is\_authenticated}, the message banners, and \code{\{\% block content \%\}}.
Every other template begins \code{\{\% extends 'base.html' \%\}}.

\begin{jargon}{Template inheritance}
A child template declares a parent and fills in named blocks the parent left open. The
navigation bar and footer are written once. This is the same idea as a base class in
object-oriented code.
\end{jargon}

\begin{jargon}{Auto-escaping}
Django templates escape HTML by default: a donor named \code{<script>alert(1)</script>}
renders as visible text, not executing script. This is what makes the dashboard safe
against stored XSS (cross-site scripting) despite displaying user-entered strings
everywhere, and it is on by default rather than by any decision in this project.
\end{jargon}

The pagination links do something worth noticing. To page through a filtered result set,
the current filters must survive the click:

\begin{lstlisting}[caption={\file{apps/donations/templates/donations/dashboard.html}}]
href="?page={{ page_obj.next_page_number }}{% for key, value in request.GET.items %}{% if key != 'page' %}&{{ key }}={{ value }}{% endif %}{% endfor %}"
\end{lstlisting}

It re-emits every current query parameter except \code{page}. It works. It is also the
kind of logic that belongs in a template tag rather than inline, and it is repeated four
times in the file, once each for First, Previous, Next and Last.

\begin{honest}{The dashboard hard-codes PKR in three places}
The two total cards and the top-pets panel all render \code{... PKR} as literal text in
the template, regardless of what currency the underlying rows actually hold. Combined
with the free-text \code{currency} field, this produces the bug in
§\ref{sec:currencybug}. The table rows, by contrast, correctly print each donation's own
\code{\{\{ donation.currency \}\}}, so a single page can show a row labelled USD sitting
underneath a total labelled PKR that includes it.
\end{honest}

\section{The mixed-currency bug}
\label{sec:currencybug}

\begin{honest}{Totals sum across every currency and label the result PKR, verified by running it}
\code{Sum('amount')} adds the \code{amount} column with no regard for the
\code{currency} column beside it, and the template then appends the literal string
``PKR''. I drove this through the test client with one 100 PKR donation and one 100 USD
donation:
\begin{lstlisting}
total_all_time (100 PKR + 100 USD): 200
rendered page contains "200.00 PKR"? True
\end{lstlisting}
The dashboard states that the rescue has received 200.00 PKR. It has received 100 PKR and
100 USD. The number is meaningless and it is presented with total confidence, which is
worse than an error message. The same applies to \code{total\_filtered} and to every row
of the top-pets panel.
\end{honest}

\begin{honest}{The README's own account of this is subtly wrong, and the truth is more dangerous}
The README says, under ``What I'd do differently'': ``Validate \code{currency} against a
choice list. It's a free-text field defaulting to PKR, so typos (`pkr', `Rs') would
silently split the totals.'' Running it shows the reverse. Nothing splits. Everything
merges:
\begin{lstlisting}
distinct currency values stored: ['PKR', 'Rs', 'USD', 'pkr']
total shown on dashboard: one number, all four summed, labelled PKR
filter currency=PKR (uses __iexact) -> matches 2 rows: catches 'pkr', misses 'Rs'
\end{lstlisting}
The totals never split by currency because they never group by currency in the first
place. \code{currency\_\_iexact} means the \emph{filter} also folds ``pkr'' into ``PKR''
correctly, so the case-typo scenario the README worries about is the one case already
handled. The actual failure is the one the README does not mention: a genuinely different
currency is added to the total as though it were rupees. Recording this because the
README will be read by an interviewer as the author's own assessment, and defending a
self-assessment that the code contradicts is worse than having no self-assessment.
\end{honest}

The correct fix is not only a choice list. Even with \code{currency} constrained to a
clean enum, \code{Sum('amount')} across mixed currencies is still wrong. The dashboard
needs to group by currency and show one total per currency, or convert at a recorded
rate. The choice list makes the grouping reliable; it does not by itself make the
arithmetic mean anything.

\section{Testing}

\subsection{What is covered, and it genuinely passes}

I created a clean virtual environment, installed the exact pinned dependencies from
\file{requirements.txt}, and ran the suite rather than trusting the README's count.

\begin{lstlisting}[caption={Verified on this machine}]
$ python manage.py test
Ran 11 tests in 12.828s

OK
\end{lstlisting}

Eleven tests, all passing, matching the README exactly. The dependency pins install
cleanly (Django 4.2.11, psycopg2-binary 2.9.11, python-dotenv 1.0.1, gunicorn 21.2.0,
whitenoise 6.6.0, dj-database-url 2.1.0).

\begin{table}[H]
\centering
\small
\begin{tabular}{@{}lll@{}}
\toprule
\textbf{Test} & \textbf{App} & \textbf{What it establishes} \\
\midrule
\code{test\_unauthorized\_user\_redirected\_to\_login} & accounts & the login gate exists \\
\code{test\_admin\_can\_login} & accounts & valid credentials redirect \\
\code{test\_admin\_list\_requires\_login} & accounts & protected before, reachable after \\
\code{test\_cannot\_delete\_self} & accounts & self-deletion guard holds \\
\code{test\_admin\_can\_access\_dashboard} & donations & dashboard renders for staff \\
\code{test\_admin\_can\_create\_donation} & donations & create writes the right row \\
\code{test\_amount\_validation} & donations & a negative amount is refused \\
\code{test\_csv\_export\_returns\_csv\_response} & donations & headers and content \\
\code{test\_search\_by\_donor\_name} & donations & search matches donor \\
\code{test\_search\_by\_pet\_name} & donations & search matches pet \\
\code{test\_filter\_by\_payment\_method} & donations & method filter runs \\
\bottomrule
\end{tabular}
\caption{The eleven tests. All eleven pass; the coverage claim is accurate.}
\end{table}

\begin{jargon}{Test database}
Django's test runner creates a throwaway database, runs every migration into it, and
destroys it afterwards. Each test method runs inside a transaction that is rolled back at
the end, so tests cannot see each other's data. This is why \code{setUp} can create the
same admin for every test without colliding.
\end{jargon}

\subsection{Where the suite is weaker than its pass rate suggests}

\begin{honest}{Every assertion is a substring match against the whole HTML page}
\code{assertContains(response, 'John Doe')} searches the entire rendered response,
navigation bar, form labels, footer and all. The project already has a scar from this:
the README records that \code{assertNotContains(response, 'Max')} failed on a page with
no donation for the pet Max, because the filter form's ``Max Amount'' label contained the
substring. The resolution was to rename the label to ``Amount To''. I confirmed the
current labels in \file{apps/donations/forms.py} are \code{'Amount From'} and
\code{'Amount To'}, so the test passes today.

The fix is worth being honest about, because it is backwards: the \emph{user interface}
was changed to suit the test's matching strategy. The test is still fragile, and the
next fixture named after any word appearing in the site's chrome revives the same
failure. Asserting on \code{response.context['page\_obj']} would test the actual
behaviour and could not collide with page text.
\end{honest}

\begin{honest}{Three tests assert less than their names claim}
\begin{itemize}
  \item \code{test\_filter\_by\_payment\_method} filters to \code{cash} and asserts only
    that ``John Doe'' (the cash donor) is present. It never asserts that Jane Smith, the
    online donor, is \emph{absent}. A completely broken filter that ignored the parameter
    and returned every row would pass this test.
  \item \code{test\_amount\_validation} posts a negative amount and asserts only that no
    row was created. It never checks the status code or that an error was shown, so a
    view that crashed with a 500 would also pass.
  \item \code{test\_csv\_export\_returns\_csv\_response} calls the export with no filter
    applied, which is the branch that works. Neither this test nor any other exercises
    the export after a filter, which is why the empty-filter bug in
    §\ref{sec:honest} survives a green suite.
\end{itemize}
\end{honest}

\begin{honest}{The two most consequential bugs are both in untested territory}
The suite passes completely and both bugs I found are real. Nothing tests the export with
a filter applied, and nothing tests a total with more than one currency present. This is
the useful lesson from the project: eleven green tests bought real confidence about
authentication and CRUD, and no confidence at all about the two features most likely to
produce a wrong number in front of a user.
\end{honest}

\section{Honest findings, collected}
\label{sec:honest}

Ordered by how much they would matter to a reviewer.

\begin{honest}{1. A zero-match filter exports the whole database (verified)}
\code{export\_csv\_view}'s \code{if filtered\_ids:} cannot tell an empty result from an
absent session key, so filtering to nothing and clicking Export CSV downloads every
record, including every donor phone number and email address. Confirmed by execution.
This is a live bug, not a design preference.
\end{honest}

\begin{honest}{2. Mixed-currency totals are summed and labelled PKR (verified)}
100 PKR plus 100 USD renders as ``200.00 PKR''. Confirmed by execution. The README's
description of this area (``typos would silently split the totals'') is the opposite of
what the code does.
\end{honest}

\begin{honest}{3. \file{db.sqlite3} is committed, containing a password hash and live sessions}
The database file is tracked in git and \file{.gitignore} does not exclude it. It was
added deliberately (commit \code{a49871e}, ``Add local dev DB from the screenshot
verification run''). Its contents, read directly:
\begin{lstlisting}
users:           1 row  -> admin@example.com, is_staff=1, is_superuser=1,
                           password = pbkdf2_sha256$600000$... (a real hash)
django_session:  2 rows -> session keys, expired 2026-07-08
donations:       0 rows
\end{lstlisting}
The exposure is limited and worth stating precisely rather than alarmingly: the email is
the placeholder from \file{.env.example}, the donations table is empty, the sessions
expired, and PBKDF2 at 600{,}000 iterations is not casually reversible. So this is not an
emergency. It is still a committed credential hash and a bad habit: the file is
build output, it will drift with every local run, it will produce meaningless binary
diffs, and the next person who commits it after entering real donor data publishes real
donor data. \file{db.sqlite3} belongs in \file{.gitignore} and the file should be removed
from tracking; if the screenshot needs data, a fixture or a seed command is the tool.
\end{honest}

\begin{honest}{4. Web-created admins are unvalidated superusers}
A one-character password is accepted (verified), and every account created through the UI
gets \code{is\_superuser = True}. The two compound: the weakest possible password on the
most privileged possible account, in an app whose whole security model is ``there is a
login''.
\end{honest}

\begin{honest}{5. No audit trail}
\code{Donation} has no foreign key to \code{User}. Nothing records who entered or deleted
a record, and deletes are hard deletes. For a donation tracker this is the most
significant \emph{structural} gap, as opposed to the bugs above.
\end{honest}

\begin{honest}{6. \texttt{SECRET\_KEY} has a working insecure default}
A deployment that forgets the variable runs happily on a key published in the repository.
\end{honest}

\begin{honest}{7. Session-based export state}
Unbounded growth (one integer per matching row, per dashboard render, rewritten to the
database on every request thanks to \code{SESSION\_SAVE\_EVERY\_REQUEST}) and staleness
by construction. The README identifies this one accurately.
\end{honest}

\begin{honest}{8. Missing indexes}
\code{donation\_date} drives the default ordering and \code{pet\_name} drives both the
search and the top-pets grouping. Neither is indexed, so both are full table scans. At
this data volume it is invisible, and it is the first thing that would bite at scale. The
README flags this.
\end{honest}

\begin{honest}{9. Smaller things, listed for completeness}
\begin{itemize}
  \item \code{UserCreationForm} is imported in \file{apps/accounts/forms.py} and never
    used. Dead import.
  \item \code{is\_admin} is defined identically in two files.
  \item \code{clean\_amount} duplicates the model's \code{MinValueValidator} and only
    changes the error wording (verified); its \code{if amount and ...} test should be
    \code{if amount is not None and ...}.
  \item \code{logout\_view} still permits GET.
  \item \code{django.contrib.admin} is installed and routed at \code{/admin/}, but
    neither app has an \file{admin.py}, so nothing is registered and the built-in admin
    shows only groups and permissions. It is dead weight, or an unfinished intention.
  \item Form errors are surfaced twice, once with raw field names in a banner.
  \item \file{.env.example} ships \code{DEBUG=True} while \file{settings.py} defaults to
    \code{False}. Both defensible, mildly contradictory.
  \item \code{SECURE\_BROWSER\_XSS\_FILTER} sets a header every modern browser ignores.
  \item Bootstrap loads from a CDN, so the UI is unstyled offline.
\end{itemize}
\end{honest}

\section{What is genuinely good here}

An honest document says this part too, and it is not padding.

\begin{itemize}
  \item \textbf{The custom user model was done on day one.} \code{AUTH\_USER\_MODEL} is
    almost impossible to change after the first migration. Getting it right at the start
    is the single highest-leverage decision in a new Django project, and the README's
    ``What I learned'' shows the author knows why.
  \item \textbf{Money is a \code{Decimal}.} Common enough to get wrong that getting it
    right deserves the credit.
  \item \textbf{\code{is not None} on the amount filters.} A precise, deliberate guard
    against a real edge case, showing the author was thinking about falsy zero.
  \item \textbf{The \code{TESTING} flag is a genuine debugging win.} Diagnosing why every
    test suddenly 301-redirects to https requires understanding that
    \code{manage.py test} does not set \code{DEBUG}. That is a non-obvious chain and the
    README reconstructs it correctly.
  \item \textbf{The top-pets fix.} Noticing that a panel disagreed with the total beside
    it, and tracing it to aggregating over the unfiltered manager, is exactly the kind of
    bug most people ship.
  \item \textbf{The self-deletion guard is enforced server-side}, not just hidden in the
    template.
  \item \textbf{The README volunteers its own weaknesses.} Most of §\ref{sec:honest}
    items 5, 7 and 8 are the author's own findings. That is rarer than it should be, and
    it is worth defending as a deliberate practice.
\end{itemize}

\section{If this project continued}

In the order a maintainer should actually take them:

\begin{enumerate}
  \item Fix the empty-filter export. It is a live data-exposure bug and a one-line
    change. Add the test that catches it.
  \item Group totals by currency, or constrain \code{currency} and then group. Stop
    printing a number that is not true.
  \item Untrack \file{db.sqlite3} and add it to \file{.gitignore}.
  \item Validate passwords in \code{AdminCreationForm}, and stop making every admin a
    superuser.
  \item Add \code{entered\_by = ForeignKey(User, on\_delete=PROTECT)} to \code{Donation}
    and stop hard-deleting.
  \item Re-derive the export queryset from the query string and delete the session
    handoff, which retires findings 1 and 7 together.
  \item Extract the filter logic into a shared helper or adopt \code{django-filter}.
  \item Make \code{SECRET\_KEY} mandatory in production.
  \item Index \code{donation\_date} and \code{pet\_name}.
  \item Split settings into base, dev and prod, retiring the \code{sys.argv} hack.
\end{enumerate}

\begin{keyidea}{The one-sentence summary}
A small, conventional, competently built Django CRUD application that gets several
non-obvious Django decisions right (custom user model up front, Decimal money, lazy
QuerySet filtering, a real reason for its test-mode flag) and carries two real bugs in
the two places nobody tested: an empty filter exports everything, and mixed currencies
are added together and labelled PKR.
\end{keyidea}

\vfill
\begin{center}
\small\itshape
Every claim in this document was checked against the source. Where a claim is marked
verified, it was established by executing the code: the test suite, the form validation
probes, the aggregation SQL, the export behaviour, and the contents of the committed
database. The single item labelled inferred is the rationale for the mixed primary-key
styles, which no comment in the source explains.
\end{center}

\end{document}
